\Article{11}
Wenden wir uns nun näher zu unserm Gegenstande selbst. Wenn für den Modulus $m=a+bi$ die Zahlen $f,f',f''$ etc. so beschaffen sind, dass sie erstlich alle nach dem Modulus $m$ unter sich incongruent sind, zweitens aber jede ganze Zahl einer von ihnen nothwendig congruent sein muss, so nennen wir den
\SourceLeaf{326}
Inbegriff der Zahlen $f,f',f''$ etc. das System der Primitivreste von $m$. Ihre Anzahl ist immer $=aa+bb$.

\Article{12}
Man kann das System der Primitivreste auf vielfache Art bilden; die einfachste ist, die Punkte innerhalb des Quadrats $0,m,(1+i)m,im$ zu wählen; dazu müssen aber noch hinzugefügt werden
\begin{enumerate}
\item[I.] der Punkt oder die Grösse $0$
\item[II.] alle Punkte auf zwei einander nicht gegenüberliegenden Grenzlinien.
\end{enumerate}
Anstatt auf einer der $4$ Grenzlinien alle Punkte zu nehmen, kann man sie auch auf mehren zugleich nehmen.

Diese Auswahl dieser Punkte auf den Grenzlinien, falls welche darauf fallen, kann auf mehrfache Art geschehen, so dass obigen Bedingungen Genüge geschieht. Am einfachsten ist die folgende Manier.

Man nehme auf der Grenzlinie $0,m$ alle Punkte zwischen $0$ und $\frac12m$ inclus. und auf der Grenzlinie $0,im$ alle Punkte von $\frac12im$ bis $im$ exclusive und auf ähnliche Art bei den beiden andern.

Man kann diese beiden Manieren so sinnlich darstellen
\begin{center}\resizebox{.85\textwidth}{!}{\input{diagrams/native326_boundaries.tex}}\end{center}

\Article{13}
Schliesst man von den Primitivpunkten aus
\begin{enumerate}
\item[I.] Bloss den Punkt $0$, wenn $a$ gerade und $b$ ungerade oder umgekehrt.
\item[II.] Die Punkte $0$ und $\frac12(1+i)m$, wenn $a$ und $b$ beide ungerade.
\item[III.] Die vier Punkte $0,\frac12m,\frac12(1+i)m,\frac12im$, wenn $a$ und $b$ beide gerade,
\end{enumerate}
so nennen wir die übrigbleibenden eigentliche Primitivpunkte, die ausgeschlossenen uneigentliche. Die Anzahl von jenen ist also
\[
\begin{array}{rrcl}
\text{in Fall}&\mathrm{I}&=&aa+bb-1\\
&\mathrm{II}&=&aa+bb-2\\
&\mathrm{III}&=&aa+bb-4
\end{array}
\]
also immer durch $4$ theilbar.

\SourceLeaf{327}
\Article{14}
Diese eigentlichen Primitivpunkte lassen sich in $4$ Classen $F,F',F'',F'''$ theilen, so dass
\[
\begin{array}{llll}
iF\equiv F'&iF'\equiv F''&iF''\equiv F'''&iF'''\equiv F\\
-F\equiv F''&-F'\equiv F'''&-F''\equiv F&-F'''\equiv F'\\
-iF\equiv F'''&-iF'\equiv F&-iF''\equiv F'&-iF'''\equiv F''
\end{array}
\]
Hiebei findet nun folgendes höchst wichtige Theorem statt.

Es sei $M$ eine Zahl, welche mit $m$ keinen Factor gemein hat. Von den Zahlen $MF$ gehören in die Classe $F$ eine Anzahl von $n$
\[
\begin{array}{cc}F'&n'\\F''&n''\\F'''&n'''\end{array}
\]
und der kleinste Rest von $n'+2n''+3n''$ nach dem Modulus $4$ sei $=N$, also $N$ einer der $4$ Zahlen $0,1,2,3$ gleich: unter dieser Voraussetzung ist $N$ \emph{unabhängig} von der Art der Vertheilung der Primitivreste in Classen. Wir nennen ihn den Decident des biquadratischen Verhältnisses der Zahl $M$ zu $m$.

\Article{15}
Die einfachste Art der Vertheilung ist allerdings folgende
\begin{center}\resizebox{.35\textwidth}{!}{\input{diagrams/native327_classes.tex}}\end{center}
Inzwischen kann in speciellen Fällen eine andere Vertheilung vortheilhafter sein.

\Article{16}
Sind $f,f',f''$ etc. die sämmtlichen Primitivreste des Modulus $m$, so ist
\SourceLeaf{328}
\[
\begin{aligned}
&S(z+\tfrac fm,\ z'+\tfrac fm,\ z''+\tfrac fm,\ \text{etc.})\\
&+S(z+\tfrac{f'}m,\ z'+\tfrac{f'}m,\ z''+\tfrac{f'}m,\ \text{etc.})\\
&+S(z+\tfrac{f''}m,\ z'+\tfrac{f''}m,\ z''+\tfrac{f''}m,\ \text{etc.})\\
&+\text{etc.}\\
&=S(mz,mz',mz'',\ \text{etc.}
\end{aligned}
\]

